\begin{abstract}
As large language models train on ever-larger web corpora, public benchmarks increasingly leak into training data, enabling models to achieve high scores through memorization rather than genuine capability. We present the first systematic, controlled measurement of \emph{finetuning resistance} across seven ML benchmarks. We fine-tune \qwen separately on four benchmarks' training data and evaluate each resulting model on all seven benchmarks, including shadow and harder-variant benchmarks (\gsmsym for \gsmk, \mmlupro for \mmlu). We quantify the portion of each benchmark's improvement attributable to memorization versus genuine capability gain. Our key finding is that no static benchmark is truly finetuning-proof, but benchmarks vary dramatically in their resistance. \gpqa and \mmlupro are effectively immune to fine-tuning on related data, while \gsmk shows significant inflation: 34\% of its fine-tuning gain reflects memorization rather than real improvement. An unexpected finding is that fine-tuning on \emph{any} benchmark---even commonsense reasoning---dramatically improves \gsmk performance, suggesting that \gsmk's low baseline scores reflect poor instruction-following rather than lack of mathematical knowledge. We rank benchmarks by resistance and identify three design principles---expert-knowledge requirements, increased answer choices, and template randomization---that together maximize finetuning resistance.
\end{abstract}
